% Section 10.2, from lectures/L10/appendix/d_embedded_constraints.md, without its exercise, which is
% the last set of exercises in Section 10.4.

\section{Constraints on Resource-Constrained Targets}
\label{c10:sec:embedded}\label{c10:app:d}
Everything built in this book has run on a general-purpose computer, where memory is effectively
unlimited and every CPU core has a hardware floating-point unit (FPU). Deploying the same kind of
network to a microcontroller or other resource-constrained target introduces constraints that never
show up on a laptop. This section surveys the main ones, and the last exercises in
\secref{c10:sec:exercises} work through two of them by hand: memory footprint and fixed-point
arithmetic.

\subsection{Memory footprint}
\label{c10:sec:memory}
Every weight and bias your network learned during training has to be stored somewhere on the target
device, typically in flash (read-only after deployment) or RAM (if the network keeps learning on
the device). The \code{std::vector<double>} used throughout this book stores each value as an
8-byte \code{double}.

Consider the small hidden and output layer pair many of the exercises built (3 nodes with 2 weights
each in the hidden layer, 1 node with 3 weights in the output layer):

\begin{console}
Hidden layer: 3 * 2 weights + 3 biases  =  9 values
Output layer: 1 * 3 weights + 1 bias    =  4 values
Total: 13 doubles * 8 bytes             =  104 bytes
\end{console}

\noindent That's negligible on any target. But the examples in this book are intentionally tiny. A
single hidden layer sized for 28\,×\,28 pixel images (784 inputs) with 128 nodes (a very ordinary
size for a first hidden layer) tells a different story:

\begin{console}
Weights: 128 * 784 = 100,352
Biases:  128
Total:   100,480 values

As double (8 bytes each): 803,840 bytes  ≈ 785 KB
As float  (4 bytes each): 401,920 bytes  ≈ 392.5 KB
As int8   (1 byte each):  100,480 bytes  ≈ 98.125 KB
\end{console}

\noindent That's \textbf{one layer}. Many microcontrollers have on the order of tens of KB of RAM;
mid-range parts often have a few hundred KB; only higher-end parts reach into the multiple-MB range.
The \code{double} version of this single layer already exceeds what a lot of embedded targets have
available for their \emph{entire} program; and a real network has several layers. This is why
embedded machine learning almost always trains at full precision on a workstation and then
\textbf{quantizes} the trained weights down to a smaller type before deployment (see
\secref{c10:sec:quantization}), rather than training or storing at \code{double} precision on the
target itself.

\subsection{Precision: floating-point vs.\ fixed-point}
\label{c10:sec:fixedpoint}
Many microcontrollers (especially smaller Cortex-M0 and M0+ parts, or 8- and 16-bit MCUs) have no
hardware FPU. Floating-point operations on those targets are emulated in software, which can be an
order of magnitude or more slower than an integer add or multiply.

\textbf{Fixed-point} arithmetic sidesteps this by representing a real number as a scaled integer
instead of a \code{float} or \code{double}. A common convention is Q$m$.$n$ notation: $m$ integer
bits, $n$ fractional bits. In \textbf{Q16.16} format, a 32-bit signed integer represents a real
number scaled by $2^{16}$ (65536):

\begin{cppcode}
using Fixed16_16 = std::int32_t;

constexpr int kFractionalBits{16};
constexpr Fixed16_16 kOne{1 << kFractionalBits};  // 65536

// Convert a floating-point value to Q16.16 fixed-point.
Fixed16_16 toFixed(const double value) noexcept
{
    return static_cast<Fixed16_16>(value * kOne);
}

// Convert a Q16.16 fixed-point value back to floating-point.
double toDouble(const Fixed16_16 value) noexcept
{
    return static_cast<double>(value) / kOne;
}
\end{cppcode}

Addition of two fixed-point values is just integer addition: no special handling needed, since both
operands are scaled by the same factor. Multiplication needs more care: multiplying two Q16.16
values naively overflows a 32-bit integer almost immediately (two values scaled by $2^{16}$ multiply
to a result scaled by $2^{32}$), so the multiply has to happen in a wider type before scaling back
down:

\begin{cppcode}
// Multiply two Q16.16 fixed-point values.
Fixed16_16 fixedMultiply(const Fixed16_16 a, const Fixed16_16 b) noexcept
{
    const auto product = static_cast<std::int64_t>(a) * static_cast<std::int64_t>(b);
    return static_cast<Fixed16_16>(product >> kFractionalBits);
}
\end{cppcode}

The trade-off: fixed-point trades dynamic range and some precision (Q16.16 can represent roughly
$\pm 32{,}767$ with a resolution of about 0.0000153) for arithmetic that's fast and deterministic on
targets without an FPU, with no software floating-point emulation involved.

\subsection{Quantization}
\label{c10:sec:quantization}
Even on targets \emph{with} an FPU, storing weights as 8-bit integers instead of 32-bit floats cuts
memory four times, and many embedded cores can execute int8 operations faster than float32 ones.
The common approach, \textbf{post-training quantization}, trains the network at full precision as
done throughout this book, then converts the trained weights to a smaller integer type for
deployment, typically alongside a per-layer scale factor (conceptually similar to the fixed-point
conversion above, but chosen per tensor rather than a single fixed format for the whole network). A
full quantization pipeline is out of scope for this book, but it's worth recognizing the idea:
\emph{train} at high precision, \emph{deploy} at low precision.

\subsection{Real-time and latency constraints}
\label{c10:sec:realtime}
A conv layer performs on the order of \code{inputSize² × kernelSize²} multiply-adds per feedforward
pass. On a laptop that's negligible; inside a control loop with a hard deadline, e.g.\ a motor
controller running at 1\,kHz, giving the entire loop (sensor read, inference, actuation) a 1\,ms
budget, it can dominate the loop's execution time. For embedded systems, the \emph{worst-case}
execution time matters far more than the average case: a network that's usually fast but
occasionally spikes in latency can violate a real-time deadline in exactly the way that doesn't
show up during casual testing.

\subsection{Dynamic memory allocation}
\label{c10:sec:heap}
The \code{Dense}, \code{Conv}, \code{MaxPool}, and \code{Flatten} implementations built in this
book, and the standalone \code{ml::ConvLayer}, \code{ml::MaxPoolLayer} and \code{ml::FlattenLayer}
structs from Chapters~\ref{c6:ch:cnn} and~\ref{c7:ch:layers}, all use \code{std::vector}, which
allocates on the heap. On many embedded targets, heap allocation is avoided or tightly controlled:
long-running systems risk heap fragmentation over time, allocation latency isn't constant-time, and
some bare-metal targets have no heap at all. A common alternative is sizing everything at compile
time with \code{std::array} (or plain C arrays) instead of \code{std::vector}, trading the
convenience of runtime-sized layers for fully static, more predictable memory use.
